% The tree factor graph of Section~\ref{sec:phylo}, drawn by hand on the eight
% taxa of the release fixture. Input by textbook.tex; a sketch and not a
% rendered figure, so it is outside the figure cache and the render cap, which
% govern generated output only. The branch lengths are those of
% tests/regression/fixtures/tree_jc/release.yaml, and the horizontal axis is
% their unit: one node's depth is the sum of the lengths above it.
\begin{figure}[htbp]
\centering
\begin{tikzpicture}[
    x=9cm, y=0.55cm,
    fgsmall/.style={fgvariable, minimum size=0.30cm, inner sep=0pt},
    blen/.style={font=\tiny, inner sep=1pt},
    taxon/.style={anchor=west, font=\footnotesize}
  ]
  % node name/depth/height
  \foreach \n/\d/\h in {%
      root/0/3.5, l/0.05/5.5, ll/0.10/6.5, lr/0.10/4.5,
      r/0.05/1.5, rl/0.10/2.5, rr/0.10/0.5}
    { \node[fgsmall] (\n) at (\d, \h) {}; }
  \foreach \n/\d/\h/\t in {%
      A/0.20/7/A, B/0.25/6/B, C/0.30/5/C, D/0.35/4/D,
      E/0.40/3/E, F/0.20/2/F, G/0.25/1/G, H/0.50/0/H}
    { \node[fgobserved, minimum size=0.34cm, inner sep=0pt] (\n) at (\d, \h) {};
      \node[fgfactor, fill=black!55] (i\n) at (\d, \h) [xshift=7mm] {};
      \draw[fglink] (\n) -- (i\n);
      \node[taxon] at (\d, \h) [xshift=10mm] {\t}; }
  % parent/child/length: a square branch, the factor at the midpoint of its
  % horizontal segment, the length above it.
  \foreach \p/\c/\len in {%
      root/l/0.05, root/r/0.05, l/ll/0.05, l/lr/0.05, r/rl/0.05, r/rr/0.05,
      ll/A/0.10, ll/B/0.15, lr/C/0.20, lr/D/0.25,
      rl/E/0.30, rl/F/0.10, rr/G/0.15, rr/H/0.40}
    { \draw[fglink] (\p) |- (\c);
      \path (\p |- \c) -- (\c)
        node[fgfactor, pos=0.5] {} node[blen, pos=0.5, above] {\len}; }
  % the root prior
  \node[fgfactor] (prior) at (0, 3.5) [xshift=-6mm] {};
  \draw[fglink] (prior) -- (root);
  \node[fgnote, anchor=east] at (0, 3.5) [xshift=-8mm] {$\psi_\rootnode = \bpi$};
  % the axis, in the unit the branch lengths are in
  \draw[semithick] (0, -1.2) -- (0.5, -1.2);
  \foreach \d in {0, 0.1, 0.2, 0.3, 0.4, 0.5}
    { \draw[semithick] (\d, -1.2) -- ++(0, -0.25)
        node[below, font=\tiny, inner sep=1.5pt] {\d}; }
  \node[font=\footnotesize, anchor=north] at (0.25, -2.0)
    {expected substitutions per site};
\end{tikzpicture}
\caption{One site of the eight-taxon release fixture as a factor graph, drawn
at the fixture's own branch lengths. Open circles are the states of the
internal nodes and shaded circles the states at the leaves; the small square
on each branch is the pairwise factor $\Pmat(\tau_v)$ of
Equation~\ref{eq:jc}, labelled with that branch's length; the filled square at
each leaf is the indicator folding in the observed character; and the square
at the root is the prior $\bpi$. Horizontal position is the sum of the branch
lengths above a node, in the unit the axis carries, so the drawing is the
topology and the lengths at once. The bipartite graph is a tree, so
Equation~\ref{eq:sum-product} is exact and its leaf-to-root pass is
Equation~\ref{eq:pruning}. Figure~\ref{fig:sim-tree} is the same fixture with
the characters this structure generates at its leaves.}
\label{fig:phylo:sketch}
\end{figure}
